\chapter{Machine elements: fasteners, springs, joints}
\label{vol08:ch03}

\epigraph{The most important machine elements are the ones that fail when they break.}{anonymous mechanical engineer}

\chapmeta{Half-life: medium-life. Archetypes: failure, scaling. Exercise target: 35. Project track: physical. Named cases: Aloha Airlines Flight 243, Hyatt Regency walkway collapse.}

\noindent
Machine elements are the standardised parts that connect, support, and transmit motion in machines: fasteners, springs, shafts, couplings, clutches. The discipline is largely catalogue-driven: the engineer specifies what the machine needs, picks a standardised part that meets the spec, verifies the part is appropriate for the load and the lifetime, and documents the choice. Errors at this layer have produced some of the most consequential mechanical failures in the historical record. The standard reference is Shigley \cite{text:shigley-mechanical-design}; for fastener-specific design, Bickford \cite{text:bickford-bolted-joints} is unmatched.

\section{Bolted joints: preload, friction, fatigue}\pathtag{core}

A \engterm{bolted joint} clamps two or more parts together using a threaded fastener. The engineering discipline of bolted joints rests on preload: the tension induced in the bolt when it is tightened. Preload, not the bolt's tensile strength, is what carries most external load and prevents joint loosening.

\subsection{Preload}

The bolt is tightened to a preload $F_i$. The clamped parts compress under this load; the bolt stretches. When an external tensile load $P$ is applied to the joint, the bolt sees a load increase from $F_i$ to $F_i + k_b/(k_b + k_m) \cdot P$, where $k_b$ is the bolt's stiffness and $k_m$ is the clamped-material's stiffness. For typical joints $k_m / k_b \approx 5$ to $10$, so the bolt sees about $10$-$15\%$ of the external load; the rest is taken by reducing the clamping force.

The engineering implication: a properly preloaded joint shields the bolt from most of the cyclic load. The fatigue life of the bolt depends on the cyclic component, not the mean load. A joint with $F_i \gg P$ has nearly constant bolt stress; fatigue life is long. A joint with $F_i \ll P$ has bolt stress that ranges from zero (joint separated) to $P$ (full load on bolt); fatigue life is short.

\subsection{Torque-tension}

The standard tightening method is to apply a specified torque. The torque-tension relationship is:
\[
T = K \cdot d \cdot F_i,
\]
where $T$ is torque, $d$ is the nominal bolt diameter, $K$ is the nut factor (typically $0.18$-$0.20$ for clean threads, $0.10$-$0.15$ with lubricant). The nut factor varies by 30\% or more in practice depending on surface condition, lubrication, and run-in.

The torque-tension method achieves preload accuracy of approximately $\pm 25\%$. More accurate methods: angle-controlled tightening (turn-of-the-nut, $\pm 5$ to $10\%$), bolt-stretch measurement ($\pm 5\%$), strain-gauge or ultrasonic measurement ($\pm 1$ to $5\%$). The engineering choice depends on the cost of underpreload (separation, fatigue) and overpreload (yielding, stripping).

\subsection{Fatigue of bolted joints}

The fatigue failure of bolted joints follows the standard Volume V fatigue analysis with two engineering adaptations:

\textbf{Stress concentration at thread root.} The fatigue strength of a bolt is much lower than the parent material's because of the thread-root stress concentration ($K_t \approx 3$ to $5$). The first thread engaged carries about 30\% of the load; subsequent threads progressively less.

\textbf{Preload as fatigue protection.} The cyclic stress range on the bolt is $\Delta \sigma = (k_b/(k_b + k_m)) \cdot \Delta P / A$. The ratio $k_b/(k_b + k_m)$ is small when the joint is stiff (compact clamped material, soft bolt); the engineer's strategy is to keep this ratio small.

\textbf{Mean stress as fatigue concern.} The mean bolt stress is approximately the preload divided by the bolt's stress area. Material fatigue limits depend on mean stress (Goodman, Soderberg lines, Volume V Chapter 9).

\subsection{Joint loosening}

A bolted joint can loosen even without bolt failure. The mechanisms:

\textbf{Self-loosening under vibration.} The classic Junker-test result (1969): transverse vibration shakes the bolt loose far more effectively than axial vibration. Most loosening in real applications is due to transverse loading.

\textbf{Embedding.} Initial settling of asperities reduces preload by approximately $10\%$ in the first few hours of service. Modern practice: torque the joint, run it briefly, retorque.

\textbf{Creep.} Long-term plastic deformation under preload, especially at elevated temperature.

\textbf{Differential thermal expansion.} The bolt and the clamped material expand differently; the preload changes with temperature.

The remedies: lockwashers (Belleville, Nord-Lock), thread-locking compounds (Loctite), prevailing-torque nuts (nylon insert, all-metal), wire-locking, double-nutting, redundant fastening, periodic retorquing.

\section{Welded joints}\pathtag{standard}

\engterm{Welding} fuses two pieces by local melting. The result is a joint with properties similar to (but not identical to) the parent material. The engineering of welded joints rests on the heat-affected zone, the residual stresses, and the inspection regime.

\subsection{Welding processes}

\textbf{Shielded metal arc welding (SMAW, stick).} The classical manual process. Versatile, cheap, but slow and skill-dependent.

\textbf{Gas metal arc welding (GMAW, MIG).} A wire electrode fed through a torch with shielding gas. Productive, automatable.

\textbf{Gas tungsten arc welding (GTAW, TIG).} A non-consumable tungsten electrode; filler added separately. High quality, low productivity.

\textbf{Submerged arc welding (SAW).} The arc is buried under a granular flux. Used for thick-section heavy welding.

\textbf{Resistance welding.} Two electrodes pass current through the parts; resistive heating melts the contact. Used in automotive body assembly.

\textbf{Friction stir welding.} A rotating tool plunges into the joint; frictional heating softens the material; the tool traverses the joint, stirring the material together. No melting. Used for aluminium structures in aerospace and rail.

\textbf{Electron-beam and laser welding.} Deep-penetration high-energy-density processes. Used for precision and for thick sections.

\subsection{Weld joint design}

The standard joint forms: butt, lap, T, corner, edge. The standard weld types: fillet, groove (square, V, U, J), plug, slot. The engineer specifies on a drawing using ISO 2553 or AWS A2.4 weld symbols.

The strength of a weld depends on: weld type and size, parent material, filler material, heat-affected zone metallurgy, residual stress, inspection regime. Modern welding-engineering practice is governed by codes (AWS D1.1 for structural steel, ASME Section IX for pressure vessels, ISO 15614 for procedure qualification).

\subsection{Welding defects}

\textbf{Porosity.} Gas trapped in the solidifying weld. Caused by contamination, inadequate shielding.

\textbf{Slag inclusions.} Solidified flux trapped in the weld.

\textbf{Lack of fusion.} Incomplete melting at the joint line. Reduces effective weld throat.

\textbf{Undercut.} A groove melted into the parent material at the weld toe. Acts as a stress raiser.

\textbf{Cracking.} Hot cracking (during solidification), cold cracking (after cooling, hydrogen-induced), reheat cracking, lamellar tearing.

\textbf{Distortion.} Residual stresses from the welding heat cycle deform the assembly.

The engineering response: select the right process, qualify the welder, qualify the procedure, inspect (visual, dye penetrant, magnetic particle, radiography, ultrasonics), heat-treat to relieve residual stress when needed.

\section{Adhesive joints}\pathtag{standard}

\engterm{Adhesive bonding} joins surfaces through a thin layer of polymer or chemical that bonds to both. Adhesives are increasingly competitive with welds and fasteners in modern engineering: aerospace structures, automotive body panels, electronics assembly.

\subsection{Adhesive families}

\textbf{Epoxies.} Two-part or one-part heat-cured. High strength, good environmental resistance. The workhorse of structural adhesives.

\textbf{Polyurethanes.} Flexible, impact-resistant. Used in automotive glass bonding.

\textbf{Cyanoacrylates (superglue).} Single-part, room-temperature cure, very fast. Limited environmental resistance.

\textbf{Anaerobic adhesives.} Cure in the absence of oxygen and presence of metal. Used as thread-lockers and pipe sealants (Loctite family).

\textbf{Silicones.} Flexible, temperature-resistant, environmentally robust. Used in sealing applications.

\subsection{Joint design for adhesives}

Adhesives carry load best in shear, less well in tension, poorly in peel. The engineering design rule is to load the joint in shear: lap joints, double-lap, scarf joints, butt joints with reinforcement.

The surface preparation is critical: cleanliness, surface roughness, surface energy (often improved by plasma or chemical treatment). The bond's strength is set by the weaker of: cohesion (failure within the adhesive), adhesion (failure at the bond line), parent-material failure.

The adhesive's environmental tolerance is the working limitation: temperature, humidity, UV, chemical exposure, fatigue. Aerospace adhesives are qualified through extensive environmental testing; consumer adhesives often are not.

\section{Springs: helical, leaf, Belleville, gas}\pathtag{standard}

A \engterm{spring} stores energy elastically: deflection is proportional to load (for linear springs); the stored energy returns when the load is removed.

\subsection{Spring types}

\textbf{Helical compression spring.} The standard. Cylindrical coil; compression load applied along the axis. Wire diameter, coil diameter, free length, and active coil count are the design parameters.

\textbf{Helical extension spring.} Cylindrical coil; tension load; has hooks or loops at the ends.

\textbf{Helical torsion spring.} Cylindrical coil with arms; loaded in torsion about the axis. Used in mousetraps, clothespins.

\textbf{Leaf spring.} A flat strip loaded as a cantilever or simple beam. Used in vehicle suspensions.

\textbf{Belleville (disc) spring.} A conical disc; very high force per deflection. Used in bolt preloading, valve trains.

\textbf{Gas spring.} A piston in a gas-charged cylinder. Used in vehicle hatch struts, office chairs.

\textbf{Elastomer spring.} A block of rubber or polyurethane in compression. Non-linear, damped.

\textbf{Volute, constant-force, clock springs.} Specialised.

\subsection{Helical-spring design}

The standard helical-compression-spring equations: the spring rate $k = G d^4 / (8 D^3 N)$, where $G$ is the shear modulus, $d$ is the wire diameter, $D$ is the mean coil diameter, $N$ is the number of active coils. The maximum shear stress at the wire is $\tau = (8 K_s F D)/(\pi d^3)$, where $K_s$ is the Wahl correction factor accounting for wire curvature.

The fatigue design follows the same logic as bolts: identify the cyclic stress range and the mean stress; compare to the material's S-N curve with appropriate stress-concentration factor. Spring steels have well-characterised fatigue properties; the design is largely a catalogue lookup against ISO 26909 or similar.

\subsection{Buckling}

A long helical compression spring can buckle. The critical slenderness is approximately $L_f / D < 4$ for safe operation against buckling (guided ends); $L_f / D < 2.6$ for unguided ends. Springs exceeding this should be guided in a tube or on a rod.

\section{Shaft design and keyways}\pathtag{standard}

A \engterm{shaft} transmits torque or rotational motion. The design problem is to size the shaft for combined bending and torsional loads, fatigue, and stiffness, while accommodating the keyways, shoulders, splines, and bearing seats.

\subsection{Shaft design under combined loading}

A typical shaft sees: torsion from the transmitted torque, bending from radial loads on attached components (gears, pulleys), and sometimes axial loads. The combined-loading stress at any cross-section is found by superposition of the individual stresses; the maximum is at the surface of the shaft.

The fatigue design uses the modified-Goodman or ASME-elliptic approach: identify the mean and alternating components of bending and torsion stress; combine; compare to the material's endurance limit modified for surface finish, size, reliability, temperature.

\subsection{Stress concentrations}

A shaft is full of stress concentrations: keyways, shoulders, fillets, retaining-ring grooves, splines, press-fit interfaces. Each has a stress-concentration factor that must be applied to the nominal stress. The engineer's discipline is to use generous fillet radii (typically $0.1 d$ or larger, where $d$ is the shaft diameter at the smaller cross-section), to minimise keyway depth, and to avoid sharp inside corners.

\subsection{Keyways and other torque-transmitting features}

\textbf{Square or rectangular key.} The standard; transmits torque between shaft and hub through a key sitting in matching slots.

\textbf{Woodruff key.} Half-moon shape; easier to install than square keys; lower load capacity.

\textbf{Splines.} Multiple keys integrated into the shaft (involute or straight-sided). Used for high-torque applications.

\textbf{Press fits and shrink fits.} The hub is sized smaller than the shaft and forced or thermally expanded to fit. Torque transmitted by friction.

\textbf{Polygon shafts.} The shaft has a non-circular cross-section; the matching hub form fits over it. Used in automotive applications.

The engineering choice is governed by torque magnitude, alignment requirements, assembly considerations, and cost. Keys are cheap and easy but introduce stress concentrations; splines are expensive but well-balanced.

\section{Couplings and clutches}\pathtag{standard}

\subsection{Couplings}

A \engterm{coupling} connects two shafts. The types:

\textbf{Rigid coupling.} Direct connection; requires very precise alignment. Used in precision machinery.

\textbf{Flexible coupling.} Accommodates angular and parallel misalignment; common types include elastomeric (rubber-bonded), jaw, disc, gear, beam.

\textbf{Universal joint.} Accommodates large angles; introduces velocity variation (cardan joint) or constant-velocity (Rzeppa, tripod) variants.

The selection depends on: the misalignment expected (angular, parallel, axial), torque rating, speed rating, backlash tolerance.

\subsection{Clutches}

A \engterm{clutch} engages and disengages the transmission of torque. The types:

\textbf{Friction clutch.} Two plates pressed together; torque transmitted by friction. The dominant automotive clutch; also used in clutch packs in automatic transmissions.

\textbf{Dog clutch.} Engages by interlocking teeth; can only engage at low relative speed.

\textbf{Centrifugal clutch.} Engages automatically above a threshold speed.

\textbf{Magnetic clutch.} Engaged by an electromagnet.

\textbf{Fluid coupling.} Uses fluid to transmit torque; slip-controlled.

\textbf{One-way (overrunning) clutch.} Transmits torque in one direction only; freewheels in the other.

The engineering choice is workload-dependent. Clutches have a design lifetime under the engagement-disengagement duty cycle; the friction material wears and must be replaced.

\section{Failure: bolt-loosening accidents, weld-failure aviation, joint creep}\pathtag{core}

\subsection{Hyatt Regency walkway collapse, 1981}

The Hyatt Regency Hotel in Kansas City suffered a walkway collapse on 17 July 1981 that killed 114 people and injured more than 200 \cite{acc:nbs-hyatt-1982}. The collapsed structure was two elevated walkways suspended from the ceiling by steel rods. The original design called for continuous rods passing through both walkways; the as-built construction used separate rods connecting the upper walkway to the ceiling and the lower walkway to the upper walkway, doubling the load on the upper walkway's connections.

The failure mode was a hanger-rod-and-nut connection: the upper walkway's box beams were not designed to carry the lower walkway's load. The investigation by the National Bureau of Standards \cite{acc:nbs-hyatt-1982} identified the design change (a shop-drawing modification, accepted by the engineer of record without a re-analysis) as the proximate cause.

The engineering lesson is multi-layer: the integrity of bolted and threaded connections depends on the load path being correctly modelled; shop drawing changes that modify the load path require re-analysis; the engineer of record's verification of shop drawings is a critical control. Volume X Chapter 5 returns to this case at greater depth.

\subsection{Aloha Airlines Flight 243, 1988}

The Aloha 737 fuselage rupture on 28 April 1988 \cite{acc:ntsb-aar94-04} resulted from fatigue cracks at the row-of-rivets joining the upper fuselage skin to the underlying structure. The skin tore loose, exposing the cabin to 24,000-foot altitude; one flight attendant was lost, but the aircraft landed.

The mechanism was multi-site fatigue damage: many of the riveted skin lap joints had small fatigue cracks, each below the inspection threshold individually, but coupled through the shared structure into a catastrophic failure when one cracked sufficiently to load the adjacent joints beyond their reduced strength.

The lesson for bolted/riveted/welded connections is that fatigue is a distributed phenomenon: a single joint's life is not the same as the joint-family's life. The aviation-industry response was widespread structural-integrity programmes, the Aircraft Structural Integrity Program (ASIP), and the Aging Aircraft Working Group. Volume X Chapter 2 takes this as a foundational case.

\subsection{Joint creep}

A bolted joint at elevated temperature exhibits creep: the clamped material deforms over time, the bolt's elongation cannot follow, the preload drops. The classical case: high-temperature bolted flanges in steam piping or in jet-engine casings. The engineering response is to specify creep-resistant materials, to size the bolts conservatively, and to retorque on scheduled intervals.

A similar pattern is seen in adhesive joints under sustained load: long-term creep can produce significant displacement even at room temperature, especially in flexible adhesives.

\begin{principle}[Connections are where machines fail]
The interface between parts is the location of the failure mode in most mechanical accidents. A bolt that loosens, a weld that cracks, an adhesive joint that creeps, a shaft that fatigues at the keyway: each is a failure at the connection between parts. The engineering discipline is to engineer the connection with the same rigour as the parts it connects, and to verify the connection's integrity over the machine's expected lifetime.
\end{principle}

\section{Project}\pathtag{core}

\begin{project}[Design a bolted joint to withstand fatigue loading]
Design a bolted joint for a specified static and cyclic load. Build it. Measure the actual preload as a function of torque.

\begin{enumerate}[leftmargin=*]
  \item Specify the joint: two metal plates clamped by a bolt and nut. Static load $1\,\si{\kilo\newton}$, cyclic component $\pm 200\,\si{\newton}$, $10^6$ cycles.
  \item Choose the bolt: grade, diameter, thread pitch. Compute the required preload.
  \item Compute the torque to achieve the preload. Use $K = 0.18$ as the nut factor.
  \item Build the joint. Use a strain gauge on the bolt or a load cell under the nut to measure actual preload.
  \item Measure preload vs.\ torque for at least $5$ tightening-and-loosening cycles. Plot. Identify the variability in the nut factor.
  \item Compute the bolt's fatigue safety factor under the cyclic load using your measured preload.
\end{enumerate}

The deliverable is the design calculations, the measurements, the variability analysis, and a short report (no more than four pages) on the engineering implications of the measured preload variability.
\end{project}

\section{Exercises}

\subsection*{Calculation}\pathtag{core}

\begin{exercise}
Compute the required preload to maintain a $5\,\si{\kilo\newton}$ external load on a joint with bolt stiffness ratio $k_b/(k_b+k_m) = 0.15$ and the joint not to separate at $1.5\times$ design load.
\end{exercise}

\begin{exercise}
For an M10 grade 8.8 bolt with $K = 0.20$, compute the tightening torque for $25\,\si{\kilo\newton}$ preload.
\end{exercise}

\begin{exercise}
A helical compression spring has wire diameter $4\,\si{\milli\meter}$, mean coil diameter $30\,\si{\milli\meter}$, $10$ active coils, $G = 80\,\si{\giga\pascal}$. Compute spring rate.
\end{exercise}

\begin{exercise}
A shaft must transmit $500\,\si{\newton\meter}$ torque with safety factor 4 against yield. Compute minimum diameter for steel with yield $400\,\si{\mega\pascal}$.
\end{exercise}

\begin{exercise}
Compute the throat thickness of a fillet weld required to carry $20\,\si{\kilo\newton}$ shear per metre length, given allowable shear stress $100\,\si{\mega\pascal}$.
\end{exercise}

\begin{exercise}
A spline has $20$ teeth, $40\,\si{\milli\meter}$ pitch diameter, tooth height $2\,\si{\milli\meter}$. Compute the torque capacity at allowable bearing stress $50\,\si{\mega\pascal}$.
\end{exercise}

\subsection*{Derivation}\pathtag{standard}

\begin{exercise}
Derive the bolt-load increase formula $F_b = F_i + (k_b/(k_b+k_m))P$ from compatibility of displacements.
\end{exercise}

\begin{exercise}
Derive the spring rate of a helical compression spring from Castigliano's theorem applied to a single coil.
\end{exercise}

\begin{exercise}
Show that the maximum stress in a circular shaft under combined bending and torsion is $\sigma_{\max} = \sqrt{\sigma^2 + 4\tau^2} / 2$ (maximum shear stress theory).
\end{exercise}

\begin{exercise}
Derive the Junker-test result: a transverse-vibration load loosens a bolt much faster than an axial load.
\end{exercise}

\begin{exercise}
Derive the buckling slenderness limit $L_f/D < 4$ for a guided-ends helical spring under axial load.
\end{exercise}

\subsection*{Estimation}\pathtag{standard}

\begin{exercise}
Estimate the preload loss due to embedding in the first 24 hours of a typical bolted joint.
\end{exercise}

\begin{exercise}
Estimate the fatigue life of an M10 grade 8.8 bolt at $25\,\si{\kilo\newton}$ preload and $\pm 2\,\si{\kilo\newton}$ cyclic external load with stiffness ratio $0.15$.
\end{exercise}

\begin{exercise}
Estimate the relative cost of a friction-stir-welded vs.\ a riveted aluminium aerospace joint.
\end{exercise}

\begin{exercise}
Estimate the preload variation across a set of $20$ bolts on a flange, given a $\pm 25\%$ torque-tension uncertainty.
\end{exercise}

\subsection*{Simulation}\pathtag{standard}

\begin{exercise}
Simulate the load distribution among threads in a bolted joint as a function of bolt-vs-nut stiffness.
\end{exercise}

\begin{exercise}
Simulate the helical-spring stress distribution; verify the Wahl factor.
\end{exercise}

\begin{exercise}
Simulate a Belleville-spring stack: stack-in-series for greater deflection, stack-in-parallel for greater force.
\end{exercise}

\begin{exercise}
Simulate the load distribution on a riveted lap joint under axial load; identify the highest-loaded rivet.
\end{exercise}

\begin{exercise}
Simulate the Junker test: transverse vibration applied to a preloaded bolted joint; observe preload loss over cycles.
\end{exercise}

\subsection*{Design}\pathtag{standard}

\begin{exercise}
Design a bolted flange joint for $50\,\si{\kilo\newton}$ axial load and $30\,\si{\kilo\newton\meter}$ moment; specify bolts, pattern, preload.
\end{exercise}

\begin{exercise}
Design a welded joint connecting two structural-steel members carrying $200\,\si{\kilo\newton}$ at $100$ cycles per day; specify weld type, size, and inspection.
\end{exercise}

\begin{exercise}
Design an automotive valve spring for $30\,\si{\milli\meter}$ stroke at $6000\,\si{rpm}$; verify against fatigue.
\end{exercise}

\begin{exercise}
Design a press-fit hub-shaft connection for $200\,\si{\newton\meter}$ torque transmission; specify interference and assembly method.
\end{exercise}

\begin{exercise}
Design a flexible coupling between an electric motor and a centrifugal pump shaft; specify the coupling type and the alignment tolerance.
\end{exercise}

\subsection*{Judgement}\pathtag{core}

\begin{exercise}
The Hyatt Regency walkway collapse \cite{acc:nbs-hyatt-1982} involved a load-path change at the shop-drawing stage. State three engineering practices that should have caught the change.
\end{exercise}

\begin{exercise}
The Aloha Airlines Flight 243 incident \cite{acc:ntsb-aar94-04} illustrated multi-site fatigue damage in riveted joints. State the engineering implications for inspection strategy.
\end{exercise}

\begin{exercise}
A team is choosing between bolted and welded for a structural connection in a vibration environment. State the engineering considerations.
\end{exercise}

\begin{exercise}
A team's bolted joints loosen periodically; the proposed remedy is to add lockwashers. State the engineering analysis and three remediation options.
\end{exercise}

\begin{exercise}
A team uses adhesive bonding for a flight-critical aerospace structure. State the engineering responsibilities at design, manufacture, and inspection.
\end{exercise}

\begin{exercise}
A team's springs fail at $10^4$ cycles instead of the expected $10^7$. State three engineering hypotheses for the discrepancy.
\end{exercise}

\begin{exercise}
A senior engineer argues that ``preload accuracy of $\pm 25\%$ is enough for any joint.'' Critique. State the joints where higher accuracy is justified.
\end{exercise}

\begin{exercise}
A team is selecting between a flexible coupling and a precision rigid coupling for a metrology spindle. State the engineering trade-offs.
\end{exercise}

\begin{exercise}
A team's joint creep at elevated temperature reduces preload by $40\%$ in one year. State three engineering responses.
\end{exercise}

\begin{exercise}
A team uses friction-stir welding for an aluminium aerospace structure to avoid the heat-affected-zone issues of fusion welding. State the engineering advantages and the remaining qualification challenges.
\end{exercise}
